\documentclass[11pt]{article}

\usepackage[margin=1in]{geometry}
\usepackage[T1]{fontenc}
\usepackage{lmodern}
\usepackage{microtype}
\usepackage{hyperref}
\usepackage{booktabs}
\usepackage{array}
\usepackage{xcolor}
\usepackage{enumitem}
\usepackage{longtable}
\usepackage{amsmath}

\hypersetup{
  colorlinks=true,
  linkcolor=blue!60!black,
  urlcolor=blue!60!black,
  citecolor=blue!60!black
}

\title{\textbf{Toward Verifiable Formula Extraction for Docling}\\
\large A Steering Document for Dataset, Adapter, and Validation Work}
\author{Frontier Applications Lab Working Draft}
\date{June 13, 2026}

\begin{document}
\maketitle

\begin{abstract}
Scientific documents contain mathematical expressions that are easy for humans to read but difficult for machines to recover faithfully. Existing document parsers and small vision-language models can already extract formulas into \LaTeX, but the field still lacks strong validation loops, hard real-world benchmarks, and semantic trust signals. This document frames a research and product direction: build a Docling-compatible formula extraction adapter whose first-class asset is not merely a model, but a curated dataset and validation harness that converts formula pixels into \LaTeX, checks renderability, measures visual equivalence, canonicalizes outputs, and selectively validates semantics with Lean where typed formalization is possible.
\end{abstract}

\section{Position}

The opportunity is to treat formula extraction as compilation from visual mathematical notation into structured representations. A useful system should produce \LaTeX, preserve document location, expose confidence and validation metadata, and support semantic validation for a subset of expressions.

\begin{quote}
\textbf{Program thesis:} build verifiable structured extraction from scientific documents, beginning with mathematical formulas and expanding to code, tables, charts, and proof-carrying technical knowledge.
\end{quote}

\section{Seed Context}

Taru's seed links point to Docling's code/formula stack: CodeFormula, CodeFormulaV2, SynthFormulaNet, SynthCodeNet, ChartNet, Segment Anything, and the 2024 Vary-toy paper on reinforced vision vocabularies.

The immediate implication is that the first serious step is dataset curation and evaluation, not training a small model from scratch. CodeFormulaV2 and SmolDocling already provide strong baselines. Our edge should come from hard slices, validation, semantic structure, and integration quality.

\section{Existing Baselines}

\begin{longtable}{p{0.22\linewidth} p{0.2\linewidth} p{0.48\linewidth}}
\toprule
\textbf{Asset} & \textbf{Type} & \textbf{Relevance} \\
\midrule
CodeFormula & Model & 0.2B-class model for 120 DPI image crops of code or formulas. Outputs code or \LaTeX. \\
CodeFormulaV2 & Model & 0.3B-class newer baseline trained on SynthFormulaNet and SynthCodeNet. Best first baseline. \\
SynthFormulaNet & Dataset & 6.45M image-to-\LaTeX rows. Core source for sampling, validation, and benchmark construction. \\
SynthCodeNet & Dataset & 9.33M image-to-code rows. Useful analogy for structured visual-code transcription. \\
SmolDocling & Model/system & 256M document VLM that emits DocTags for full-page document conversion. Strategic integration anchor. \\
ChartNet & Dataset & 1.5M chart dataset with image, code, data table, summary, and QA. Strong dataset-design template. \\
\bottomrule
\end{longtable}

\section{Docling Adapter Surface}

Docling exposes code/formula enrichment through pipeline options. A practical integration target is:

\begin{verbatim}
from docling.datamodel.pipeline_options import CodeFormulaVlmOptions
from docling.datamodel.pipeline_options import PdfPipelineOptions

code_formula_options = CodeFormulaVlmOptions.from_preset("codeformulav2")

pipeline_options = PdfPipelineOptions(
    do_ocr=False,
    do_code_enrichment=True,
    do_formula_enrichment=True,
    code_formula_options=code_formula_options,
)
\end{verbatim}

The adapter should return formula text into Docling's document representation while retaining provenance, bounding boxes, confidence, validation status, and possibly alternate candidates.

\section{Validation Stack}

\subsection{LaTeX Validation}

The first validation layer checks whether predicted \LaTeX compiles or renders under a controlled engine. This catches broken commands, unbalanced braces, unsupported environments, and malformed structures.

\subsection{Visual Equivalence}

\LaTeX exact match is too brittle. Equivalent strings may render to the same expression. The validation harness should render canonical \LaTeX and compare it to the original crop using image similarity and layout-aware metrics.

\subsection{Canonicalization}

Canonicalization should normalize whitespace, macro aliases, bracket sizing, optional tags, and common equivalent constructs. This is an engineering asset: it improves both training data quality and evaluation stability.

\subsection{Lean Validation}

Lean validates typed formal mathematics, not \LaTeX. It should be used on a curated semantic subset where variables, domains, and assumptions can be recovered or provided.

\begin{center}
\begin{tabular}{lll}
\toprule
\textbf{Layer} & \textbf{Coverage} & \textbf{Trust signal} \\
\midrule
\LaTeX compile & Broad & Syntactic/renderability trust \\
Visual match & Broad & Perceptual transcription trust \\
Canonicalization & Broad & Evaluation stability \\
Lean check & Narrow & Typed semantic trust \\
\bottomrule
\end{tabular}
\end{center}

\section{Dataset Schema}

Each row should support reproducibility, validation, and future model training:

\begin{verbatim}
id
source_dataset
source_license
image_path_or_bytes
raw_latex
canonical_latex
docling_target_text
bbox_tokens
dpi
render_engine
render_status
latex_compile_status
visual_similarity_score
formula_type
complexity_tags
lean_candidate
lean_statement
lean_check_status
split_group
notes
\end{verbatim}

\section{Opportunity Matrix}

\begin{longtable}{p{0.22\linewidth} p{0.2\linewidth} p{0.18\linewidth} p{0.28\linewidth}}
\toprule
\textbf{Workstream} & \textbf{Value} & \textbf{Difficulty} & \textbf{Notes} \\
\midrule
Curated formula benchmark & High & Medium & Best first public artifact. Shows taste and creates leverage. \\
LaTeX render validator & High & Low-Medium & Required before model claims are credible. \\
Visual equivalence scoring & High & Medium & Harder than string metrics; useful for robust evaluation. \\
Docling adapter wrapper & High & Medium & Converts research into usable infrastructure. \\
CodeFormulaV2 baseline report & High & Low & Establishes known performance and failure slices. \\
Targeted synthetic generator & High & Medium & Generate hard formulas after observing failures. \\
Lean-valid subset & Very high & High & Narrow coverage but strategically distinctive. \\
Fine-tuned small adapter & Medium-High & Medium-High & Valuable only after dataset/eval exposes specific gaps. \\
Tool-interleaved recognizer & Medium-High & High & Inspired by DianJin-OCR-R1; use experts to cross-check predictions. \\
Full scientific document knowledge graph & Very high & Very high & Later-stage expansion after formula extraction works. \\
\bottomrule
\end{longtable}

\section{Strategic Bets}

\subsection{Bet 1: Dataset Quality Beats Another Model Demo}

The field already has many model demos. A benchmark with hard, validated, domain-stratified formulas can become a durable asset.

\subsection{Bet 2: Formula Extraction Needs Context}

Isolated crops are useful but incomplete. Many formulas depend on page context: variable declarations, assumptions, equation numbering, captions, and surrounding prose. The adapter should eventually support both crop-only and context-aware modes.

\subsection{Bet 3: Lean Is A Differentiator, Not A Universal Filter}

Lean validation is powerful precisely because it is strict. Use it where the semantics can be made explicit. Keep it separate from LaTeX validation so the project does not collapse under impossible coverage expectations.

\subsection{Bet 4: Compression And Reading Order Are The Frontier}

The latest descendant papers around Vary-toy point toward efficient visual tokens, token pruning, optical compression, and semantic token reordering. Formula recognition is a natural testbed because mathematical layout is inherently two-dimensional.

\section{Roadmap}

\subsection{Phase 0: Ground Truth Inspection}

Sample 1,000--10,000 SynthFormulaNet rows. Manually inspect formula diversity, target formatting, location tokens, rendering consistency, and mismatch cases.

\subsection{Phase 1: Validation Harness}

Build a local pipeline that:

\begin{enumerate}[leftmargin=*]
\item parses dataset rows,
\item extracts raw and canonical \LaTeX,
\item renders formula images,
\item records compile/render failures,
\item computes basic image similarity,
\item stratifies formula complexity.
\end{enumerate}

\subsection{Phase 2: Baseline Evaluation}

Run CodeFormulaV2 and at least one larger general VLM baseline on a small curated split. Score exact match, normalized edit distance, compile rate, render similarity, and failure tags.

\subsection{Phase 3: Hard-slice Dataset}

Create a benchmark with real and synthetic examples:

\begin{itemize}[leftmargin=*]
\item nested fractions,
\item matrices and arrays,
\item aligned equations,
\item superscript/subscript nesting,
\item integrals, sums, products,
\item piecewise definitions,
\item Greek/Latin ambiguity,
\item noisy scanned crops,
\item formulas with nearby equation numbers,
\item theorem statements with assumptions.
\end{itemize}

\subsection{Phase 4: Adapter Prototype}

Wrap the best recognizer as a Docling-compatible formula enrichment stage with validation metadata. If needed, LoRA/adapt on hard slices rather than the full dataset.

\subsection{Phase 5: Lean Subset}

Select a small formalizable subset and build a path from formula to Lean:

\[
\text{image} \rightarrow \text{LaTeX} \rightarrow \text{canonical form} \rightarrow \text{typed semantic parse} \rightarrow \text{Lean statement}
\]

This phase should begin with elementary algebra/calculus identities and theorem-like expressions where types and assumptions are explicit.

\section{Evaluation}

\begin{longtable}{p{0.28\linewidth} p{0.22\linewidth} p{0.4\linewidth}}
\toprule
\textbf{Metric} & \textbf{Scope} & \textbf{Purpose} \\
\midrule
Raw exact match & All examples & Simple but too strict. \\
Canonical exact match & All examples & Better string metric after normalization. \\
Token edit distance & All examples & Measures closeness of transcription. \\
\LaTeX compile rate & All examples & Detects malformed output. \\
Render similarity & All examples & Captures visual equivalence. \\
Complexity-slice accuracy & All examples & Reveals hidden failure modes. \\
Lean type-check rate & Semantic subset & Measures formal validity where applicable. \\
Docling integration success & PDFs & Measures practical system behavior. \\
\bottomrule
\end{longtable}

\section{What We Should Not Do First}

\begin{itemize}[leftmargin=*]
\item Do not train a model from scratch before building the evaluation harness.
\item Do not claim Lean validates arbitrary \LaTeX.
\item Do not treat synthetic data scale as equivalent to real-world robustness.
\item Do not build a generic PDF parser when Docling already provides the document pipeline.
\item Do not optimize for leaderboard metrics before building failure visibility.
\end{itemize}

\section{Near-term Deliverables}

\begin{enumerate}[leftmargin=*]
\item A reproducible dataset sampler for SynthFormulaNet.
\item A \LaTeX compile/render validator.
\item A 500--2,000 example hard benchmark.
\item A CodeFormulaV2 baseline report.
\item A failure taxonomy.
\item A Docling formula enrichment adapter prototype.
\item A small Lean-valid pilot subset.
\end{enumerate}

\section{Open Questions}

\begin{itemize}[leftmargin=*]
\item Which formula classes matter most for the target users: academic PDFs, textbooks, engineering docs, finance reports, or handwritten notes?
\item Should the first public artifact be a benchmark, an adapter, a validation tool, or a report?
\item Can visual equivalence be reliable enough without human review?
\item How much context is needed to make Lean validation useful?
\item Should the model emit one best \LaTeX string or multiple candidates with validator-guided reranking?
\end{itemize}

\section{Source Notes}

Core sources include CodeFormula, CodeFormulaV2, SynthFormulaNet, SynthCodeNet, SmolDocling, Vary-toy, GOT, DeepSeek-OCR, DeepSeek-OCR 2, DocTron-Formula, ChartNet, and the Docling code/formula enrichment documentation.

\begin{thebibliography}{99}

\bibitem{codeformula}
Docling Project. CodeFormula model card. \url{https://huggingface.co/docling-project/CodeFormula}

\bibitem{codeformulav2}
Docling Project. CodeFormulaV2 model card. \url{https://huggingface.co/docling-project/CodeFormulaV2}

\bibitem{synthformula}
Docling Project. SynthFormulaNet dataset card. \url{https://huggingface.co/datasets/docling-project/SynthFormulaNet}

\bibitem{synthcode}
Docling Project. SynthCodeNet dataset card. \url{https://huggingface.co/datasets/docling-project/SynthCodeNet}

\bibitem{docling-code-formula}
Docling documentation. Code and formula enrichment example. \url{https://docling-project.github.io/docling/examples/code_formula_granite_docling/}

\bibitem{varytoy}
Wei et al. Small Language Model Meets with Reinforced Vision Vocabulary. arXiv:2401.12503. \url{https://arxiv.org/abs/2401.12503}

\bibitem{sam}
Kirillov et al. Segment Anything. arXiv:2304.02643. \url{https://huggingface.co/papers/2304.02643}

\bibitem{chartnet}
ChartNet: A Million-Scale, High-Quality Multimodal Dataset for Robust Chart Understanding. arXiv:2603.27064. \url{https://arxiv.org/abs/2603.27064}

\bibitem{smoldocling}
Nassar et al. SmolDocling: An ultra-compact vision-language model for end-to-end multi-modal document conversion. arXiv:2503.11576. \url{https://arxiv.org/abs/2503.11576}

\bibitem{got}
Wei et al. General OCR Theory: Towards OCR-2.0 via a Unified End-to-end Model. arXiv:2409.01704. \url{https://arxiv.org/abs/2409.01704}

\bibitem{deepseekocr}
Wei, Sun, and Li. DeepSeek-OCR: Contexts Optical Compression. arXiv:2510.18234. \url{https://arxiv.org/abs/2510.18234}

\bibitem{deepseekocr2}
Wei, Sun, and Li. DeepSeek-OCR 2: Visual Causal Flow. arXiv:2601.20552. \url{https://arxiv.org/abs/2601.20552}

\bibitem{doctronformula}
Zhong et al. DocTron-Formula: Generalized Formula Recognition in Complex and Structured Scenarios. arXiv:2508.00311. \url{https://arxiv.org/abs/2508.00311}

\bibitem{mocr}
Zheng et al. Multimodal OCR: Parse Anything from Documents. arXiv:2603.13032. \url{https://arxiv.org/abs/2603.13032}

\bibitem{lean}
Lean FRO. Learn Lean. \url{https://lean-lang.org/learn/}

\end{thebibliography}

\end{document}
